\section{Fourier formulas for the deflated curvature}
\label{gt-fourier-curvature}

\subsection{The complete two-point quantity}

On a counted pair disk, let \(F\) be the field on one slice,
\(A=F/Q\), and \(B=A'/A\) as in
Section~\ref{gt-source-curvature}. This notation also applies to
a finite Fourier readout with its own counted disk. For an analytic
function \(f\), put \(\mathcal L[f]=(f')^2-ff''\). Then
\begin{equation}
 -K_F=\int_0^1
 \frac{\mathcal L[A]}{A^2}\bigl(z_2+s(z_1-z_2)\bigr)\,\dd s .
 \label{gt-deflated-curvature}
\end{equation}
Indeed \(-B'=\mathcal L[A]/A^2\), and the complex segment between
the pair lies inside the zero-free disk for \(A\). Integrating the
derivative along the segment gives the divided difference, and
the double case gives the value at \(c\). For a distinct real pair
the integral is the real-segment average. When \(z_1,z_2=c\pm\ii b\)
with \(b>0\), conjugation instead gives
\begin{equation}
 -K_F=-\frac{\Im B(c+\ii b)}{b}
 =\frac1{2b}\int_{-b}^{b}
  \Re\!\left(\frac{\mathcal L[A](c+\ii y)}{A(c+\ii y)^2}\right)\dd y .
 \label{gt-vertical-curvature}
\end{equation}
The normalized expression on the real axis does not bound this
vertical average without another estimate.

The logarithmic derivative of \(A=F/Q\) gives
\(-B'=\mathcal L[F]/F^2-\mathcal L[Q]/Q^2\) off the zeros.
Multiplying through yields the quotient identity
\begin{equation}
 \mathcal L[A]=\frac{Q^2\mathcal L[F]-F^2\mathcal L[Q]}{Q^4},
 \qquad \mathcal L[Q]=2((z-c)^2+D).
 \label{gt-curvature-quotient}
\end{equation}
It extends across the pair by analyticity of \(A\).
The subtraction is part of the complete quantity. Estimating
its two terms separately can lose the cancellation needed at
the pair.

\subsection{A Fourier numerator at a stipulated real double}

Suppose \(F(c)=F'(c)=0,\ F''(c)\ne0\), with \(c\in\R\).
Here \(Q(z)=(z-c)^2\). Taylor expansion gives
\(A(c)=F''(c)/2,\ A'(c)=F'''(c)/6,\ A''(c)=F''''(c)/12\).
Substituting in \eqref{gt-deflated-curvature} yields
\begin{equation}
 -K_F=\frac{\mathcal P_F(c)}{18 F''(c)^2},\qquad
 \mathcal P_F(z)=2(F'''(z))^2-3F''(z)F''''(z).
 \label{gt-fourier-numerator}
\end{equation}
Thus a quantitative threshold \(-K_F\ge\kappa>0\) at a double
requires \(\mathcal P_F(c)\ge18\kappa F''(c)^2\).
A sign of the numerator alone gives only the sign of the curvature.

The general Fourier representations of Laguerre expressions
are classical \cite[Theorems~2.4 and~3.7]{csordas}. Here the
coefficients and normalization can be obtained directly.
For the finite readout \eqref{eq:gauss-readout}, suppress the
index \(N\) on distinct nodes \(x_j\) and positive weights \(\omega_j\).
Differentiating \(G_N(z)=\sum_j\omega_j\e^{\ii z x_j}\) termwise,
the \(2(G_N''')^2\) part has coefficient \(-2x_j^3x_k^3\)
and the \(-3G_N''G_N''''\) part has coefficient \(3x_j^2x_k^4\).
Averaging the coefficient with the expression interchanging \(j,k\) gives
\begin{equation}
 \mathcal P_{G_N}(z)=\frac12\sum_{j,k}\omega_j\omega_k
 x_j^2x_k^2(3x_j^2-4x_jx_k+3x_k^2)\e^{\ii z(x_j+x_k)} .
 \label{gt-atomic-numerator}
\end{equation}
Since \(3u^2-4uv+3v^2=u^2+v^2+2(u-v)^2\), the coefficients in
this sum are nonnegative. This gives nonnegative weights on the
sum frequencies. The phases themselves control the value of the sum
at a real observation.

For an even node measure, group all matching sums and then opposite
ones. For real \(x\),
\[
 \mathcal P_{G_N}(x)=C_0+2\sum_{\nu>0}C_\nu\cos(\nu x),\qquad
 C_\nu\ge0,\qquad C_0=5\sum_j\omega_j^2 x_j^6 .
\]
In fact the zero sum occurs when \(x_k=-x_j\). The matching weight
is \(\omega_k=\omega_j\), and its coefficient in
\eqref{gt-atomic-numerator} is \(5x_j^6\).
The sum over all \(j\) keeps the ordered-pair convention. The zero node,
if present, contributes zero. The remaining modes can interfere, so
this constant term does not control their total signed contribution.

For the complete normalized field \(F=F_t\) in
\eqref{eq:normalized-field}, the same differentiation and
symmetrization gives
\begin{equation}
 \mathcal P_F(z)=\frac12\iint u^2v^2(3u^2-4uv+3v^2)
                 \e^{\ii z(u+v)}\,\dd\mu_t(u)\,\dd\mu_t(v).
 \label{gt-complete-fourier-numerator}
\end{equation}
The polynomial factors and the exponential on compact sets are
integrable against the product by \eqref{eq:tail}. This justifies
all derivatives and Fubini exchanges. The pushforward under
\((u,v)\mapsto u+v\) of this positively weighted measure has no atom
at zero for a density source: \(u+v=0\) has product-measure zero.
In particular the discrete \(C_0\) is not a positive constant atom
in \eqref{gt-complete-fourier-numerator}. The complete formula
requires an oscillatory lower bound without that discrete input.

\subsection{Separate Fourier allowances and their exact limits}

The positive frequency weights suggest retaining favorable terms and
bounding the others separately. Before using such an allowance in a gap
argument, its size can be determined exactly. Write
\(m_j=\int u^j\,\dd\mu_t(u)\) in this section; these are source moments,
not the exterior reciprocal terms used in Section~\ref{sec:gt-interaction}.
For an even finite source, \eqref{gt-atomic-numerator} gives
\(\mathcal P_F(0)=3m_2m_4\). The allowance that retains the constant
frequency and pays every other cosine by \(-1\) is therefore
\begin{equation}
 \mathcal P_F(x)\ge L_{\rm abs}
 =2C_0-\mathcal P_F(0)
 =10\sum_j\omega_j^2x_j^6-3m_2m_4 .
 \label{gt-absolute-fourier-allowance}
\end{equation}
This is the value of the specified allowance, not the infimum of the
trigonometric polynomial: the phases need not attain their separate
minima simultaneously. If \(m_2m_4>0\), then \(m_6>0\), and
\[
 \omega_{\max}<\frac{3m_2m_4}{10m_6}
 \quad\Longrightarrow\quad
 L_{\rm abs}\le10\omega_{\max}m_6-3m_2m_4<0.
\]
Gaussian quadrature with at least four nodes matches the moments through
degree six, so these moments may then be evaluated in the original source.
No rate for \(\omega_{\max}\), or successful quadrature degree, follows
from this finite failure test.

For a nondegenerate even absolutely continuous source with a finite sixth
moment, the zero-frequency atom in
\eqref{gt-complete-fourier-numerator} vanishes. The same constant-mode
strategy consequently has the strictly negative allowance
\begin{equation}
 L_{\rm abs}^{\rm cont}=-3m_2m_4<0.
 \label{gt-continuous-absolute-allowance}
\end{equation}
This applies to the complete canonical source on every slice of the
closed time slab. It does not exclude a jointly estimated signed
remainder or a nontrivial neighborhood of the opposite-frequency line.
The next theorem treats the larger class that retains every favorable
phase while still paying the discarded phases separately \cite{gtphase}.

\subsection{The optimal favorable-phase allowance for the complete source}

Write \(f_t(u)=W_t(u)/Z_t\), \(a_j(u)=u^j f_t(u)\). Pushing forward the
positive weight in \eqref{gt-complete-fourier-numerator} gives an even
nonnegative density
\begin{equation}
 \varpi_t=3(a_4*a_2)-2(a_3*a_3),\qquad
 \mathcal P_{F_t}(c)=\int_\R\varpi_t(s)\cos(cs)\,\dd s,
 \qquad \mathcal W_t:=\int_\R\varpi_t=3m_2m_4.
 \label{gt-fourier-product-density}
\end{equation}
Nonnegativity follows from the joint weight before the convolution is
split; it is not inferred from the difference of two convolutions.
The source estimates also give the uniform regularity needed for
oscillatory averaging. The polynomial recurrence for the differentiated
summands of \(\Phi\) in Section~\ref{sec:source} yields, for every fixed
integer \(j\ge0\),
\[
 |\Phi^{(j)}(u)|\le C_j\exp\bigl((9+4j)|u|-\pi\e^{4|u|}\bigr).
\]
Indeed, the derivative polynomial has degree at most \(j+2\), and its
remaining sum is dominated by
\(\sum_{m\ge1}m^{2j+4}\e^{-\pi(m^2-1)}<\infty\).
The reflection identity \(p(-u)-p(u)=\sinh u\), after application of
\(\partial_u^2-1\), identifies the smooth even extension at zero.
All polynomially weighted \(L^1\) and \(L^\infty\) norms of fixed
spatial derivatives of \(a_j\) are therefore bounded uniformly for
\(0\le t\le1/2\). Here \(Z_0\le Z_t\le Z_{1/2}\), with \(Z_0>0\).
Convolution differentiation gives
\begin{align*}
 \|\varpi_t'\|_1
 &\le3\|a_4'\|_1\|a_2\|_1+2\|a_3'\|_1\|a_3\|_1,\\
 \|\varpi_t''\|_\infty
 &\le3\|a_4''\|_\infty\|a_2\|_1
       +2\|a_3''\|_\infty\|a_3\|_1.
\end{align*}
Both right sides are uniformly finite, and \(\varpi_t(s)\to0\) at either
infinity. There is a positive uniform mass floor
\begin{equation}
 \mathcal W_t\ge w_*:=
 \frac{3\bigl(\int u^2\Phi(u)\,\dd u\bigr)
          \bigl(\int u^4\Phi(u)\,\dd u\bigr)}{Z_{1/2}^{2}}>0.
 \label{gt-fourier-mass-floor}
\end{equation}

\begin{theorem}[Optimal favorable-phase allowance]
\label{gt-favorable-phase-theorem}
For \(c\ge1\), let \(E\) range over measurable subsets of
\(\{s:\cos(cs)\ge0\}\), and define
\[
 L_E(t,c)=-\mathcal W_t+
       \int_E(1+\cos(cs))\varpi_t(s)\,\dd s.
\]
Then \(L_E(t,c)\le\mathcal P_{F_t}(c)\), and uniformly on the closed
physical-time slab,
\begin{equation}
 \sup_E L_E(t,c)
 =\left(\frac1\pi-\frac12\right)\mathcal W_t+O(c^{-1}).
 \label{gt-optimal-phase-limit}
\end{equation}
In particular this best allowance is strictly negative for all
sufficiently large \(c\), with an onset independent of \(t\).
\end{theorem}
\begin{proof}
The discarded cosine is bounded below by \(-1\), giving the first
assertion. Since the retained addition is nonnegative on its permitted
domain, the largest allowance takes the entire favorable set. Let the
\(2\pi\)-periodic function \(g\) equal \(\cos\theta\) where that cosine
is nonnegative, and \(-1\) elsewhere. Its mean is
\(\eta=1/\pi-1/2<0\): the positive cosine has integral two over a
period and the complementary domain has length \(\pi\).
The primitive
\(J(\theta)=\int_0^\theta(g(v)-\eta)\,\dd v\) is bounded and periodic.
Integration by parts, using the decay and derivative bound above, gives
\[
 \sup_E L_E(t,c)-\eta\mathcal W_t
 =-\frac1c\int_\R\varpi_t'(s)J(cs)\,\dd s.
\]
Its modulus is at most
\(c^{-1}\|J\|_\infty\sup_t\|\varpi_t'\|_1\). Together with
\eqref{gt-fourier-mass-floor}, this proves the uniform conclusion.
\end{proof}

The opposite-frequency band has a more precise expansion. For fixed
\(0<\alpha\le\pi/2\), take \(E=[-\alpha/c,\alpha/c]\). Evenness gives
\(\varpi_t'(0)=0\), while direct evaluation in
\eqref{gt-fourier-product-density} gives
\[
 \varpi_t(0)=5\int_\R u^6 f_t(u)^2\,\dd u>0.
\]
With \(x=cs\), Taylor's theorem and the uniform second derivative yield
\begin{equation}
 L_E(t,c)=-\mathcal W_t+
   \frac{2(\alpha+\sin\alpha)\varpi_t(0)}{c}+O_\alpha(c^{-3}).
 \label{gt-antidiagonal-band}
\end{equation}
More explicitly the error is bounded by
\(\frac{1}{2c^3}\sup_t\|\varpi_t''\|_\infty
\int_{-\alpha}^{\alpha}x^2(1+\cos x)\,\dd x\).
Thus a positive density on the opposite-frequency line does not create
a positive atom there, and this shrinking band captures only order
\(c^{-1}\) mass.

The same limitation survives every nonnegative normalized average of
these allowances supported on \(t\in[0,1/2]\), \(c\ge c_0\): its value
is at most \(\eta w_*+C/c_0\), even for jointly measurable choices of
\(E\) depending on the center. This statement concerns lower allowances,
not the sign of \(\mathcal P_{F_t}\) itself. At a stipulated double,
division by the positive denominator in \eqref{gt-fourier-numerator}
cannot make a negative allowance certify positive curvature. A double
cannot persist over an interval because \(D'=2\) at \(D=0\).
Consequently averaging double-slice allowances is not a bound on the
full two-point drift or its time-integrated transport cost. Joint use
of the negative-phase information remains outside this exclusion.

\subsection{A negative deflated numerator with positive atomic weights}

Consider the exact even readout
\[
 G(z)=\frac{5+6\cos z+\cos 2z}{12}
     =\frac{(1+\cos z)(2+\cos z)}6.
\]
The probability masses are \(5/12\) at zero, \(1/4\) at each
of \(\pm1\), and \(1/24\) at each of \(\pm2\).
Polynomials of degrees up to four have positive norm on these five
distinct nodes. Gram--Schmidt and evenness therefore give its
irreducible five-site, zero-diagonal Jacobi matrix.
At \(c=\pi\) differentiation gives
\[
 G=G'=0,\quad G''=1/6,\quad G'''=0,\quad G''''=5/6 .
\]
This is an actual double with \(\mathcal P_G(c)=-5/12\) and
\(-K_G=-5/6\) by \eqref{gt-fourier-numerator}. The slice embeds
in an exact physical heat field by multiplying the mass at
node \(x_j\) by \(\e^{(t-T)x_j^2}\) for a fixed reference time \(T\),
without a time-dependent normalization. Its factor \(2+\cos z\)
has zeros of imaginary height \(\arcosh(2)>1\), since \(\cosh(1)<2\).
Thus the five-atom control falls outside the strip hypothesis
of Corollary~\ref{gt-fixed-strip-tube}. It rules out the deflated
sign implication from positive spectral weights, symmetry and
a real double alone.

\subsection{Positive difference kernels and positive Fourier blocks}

Let \(S=\operatorname{diag}(1,-1,1,-1,\ldots)\) on the finite
chain. Then \(SJ_NS=-J_N,\ Se_0=e_0\).
Spectral projections for opposite nodes receive equal weights.
Consequently
\[
 G_N(z)=\omega_0+2\sum_{u>0}\omega_u\cos(uz),\qquad
 1-G_N(x)=4\sum_{u>0}\omega_u\sin^2(ux/2)\ge0
 \quad (x\in\R).
\]
Here \(\omega_0=0\) when the node zero is absent. At spectral parameter
zero the characteristic polynomial satisfies
\(P_n(0)=-\beta_{n-1}P_{n-2}(0)\) with \(P_0(0)=1,\ P_1(0)=0\).
Since the links are positive, an even-dimensional chain has no zero
node. These positive blocks give the bound \(G_N(x)\le1\).

The difference kernel is also positive semidefinite. For real sample
points \(y_r\) and complex coefficients \(\alpha_r\), its form is
\[
 \sum_{r,s}\overline{\alpha_r}\alpha_s G_N(y_r-y_s)
    =\sum_j\omega_j\left|
          \sum_r\alpha_r\e^{-\ii y_r x_j}\right|^2\ge0,
\]
where \(x_j\) are the spectral nodes. The positivity here arises from
difference frequencies and the modulus square. The sum-frequency coefficients in
\eqref{gt-atomic-numerator} do not supply this pointwise square or a
positive semidefinite form for the sum-frequency expression.

There is also a precise obstruction to factoring the most immediate
coefficient matrix for the deflated numerator. Replacing the second
spectral node in \eqref{gt-atomic-numerator} by its negative gives
\[
 \mathcal P_F(x)=v(x)^*\mathsf C v(x),\qquad
 v_j(x)=\e^{-\ii x x_j},\qquad
 \mathsf C_{jk}=\frac{\omega_j\omega_kx_j^2x_k^2}{2}
                  (3x_j^2+4x_jx_k+3x_k^2).
\]
For distinct positive nodes \(p,q\), its principal determinant is
\begin{equation}
 \det\mathsf C_{\{p,q\}}
 =-\frac34\omega_p^2\omega_q^2p^4q^4
       (p-q)^2(3p^2+14pq+3q^2)<0.
 \label{gt-natural-gram-obstruction}
\end{equation}
To verify the factorization, remove the positive factor
\(\omega_p^2\omega_q^2p^4q^4/4\); the remaining difference is
\(100p^2q^2-(3p^2+4pq+3q^2)^2\).
Thus this unchanged matrix has no Gram factorization. The vectors
\(v(x)\) trace a restricted curve, so matrix indefiniteness does not
prove a negative value of the numerator on that curve. Coincident
frequencies may also admit another coefficient matrix for the same
function. Exact regrouping and correlated phase estimates are not
excluded by \eqref{gt-natural-gram-obstruction}.

A nonzero real finite exponential polynomial with zero constant
coefficient cannot be nonnegative on the whole line. Indeed, if
\(P\ge0\) then \(P^2\le\|P\|_\infty P\). The long symmetric interval
mean of every nonzero frequency is zero, whereas the mean of
\(P^2=|P|^2\) is \(\sum_\nu|c_\nu|^2>0\) after equal frequencies
have been combined. This argument also applies to incommensurable
frequencies. It excludes global positivity of a literal block
without a constant frequency. It leaves open local signed estimates.

Fej\'er's identity illustrates the additional relation among
coefficients furnished by a Fourier square:
\[
 \frac1m\left|\sum_{r=0}^{m-1}\e^{\ii r\theta}\right|^2
   =1+2\sum_{k=1}^{m-1}(1-k/m)\cos(k\theta)\ge0.
\]
Each difference \(k\) occurs \(m-|k|\) times in the product, giving
the stated weights \cite[Eq.~1.15.15]{dlmf}. Arbitrary node
frequencies do not supply this lattice or these weights.
A square with general frequencies still requires its difference
support and autocorrelation coefficients to match the target.
Binning frequencies would change the expression and incur an error.

\subsection{The undeflated expression and the remaining estimate}

Differentiating \(G_N\) twice and symmetrizing yields another
nonnegative-coefficient sum:
\[
 \mathcal L[G_N](z)=\tfrac12\sum_{j,k}\omega_j\omega_k
           (x_j-x_k)^2\e^{\ii z(x_j+x_k)} .
\]
At a double this expression vanishes because \(G_N=G_N'=0\).
Even strict pointwise positivity of this undeflated expression
can coexist with nonreal zeros. Take the four-atom positive readout
\(G(z)=(9\cos z+\cos 2z)/10\), with masses \(9/20\) at each
of \(\pm1\) and \(1/20\) at each of \(\pm2\).
For \(x\) real, \((10G')^2-(10G)(10G'')\) equals
\[
 85+36\sin x\sin 2x+45\cos x\cos 2x
  =85+(81/2)\cos x+(9/2)\cos 3x .
\]
Thus \(\mathcal L[G](x)\ge(85-81/2-9/2)/100=2/5\).
Yet the equation \(2\cos^2 z+9\cos z-1=0\) includes
\(\cos z=(-9-\sqrt{89})/4<-1\), which is attained at nonreal \(z\).

These controls delimit the finite positive-weight and real-axis
Laguerre tests. The Fourier task is to bound the complete joint
numerator relative to \(F''(c)^2\) at a double or to bound the
two-point expression across the slab. In a finite readout the
counted geometry, full coordinate transfer and physical model
defect must also be paid. The actual heat-gap law
\eqref{eq:gt-gap-law} applies to the physical field, and an
estimate on one double slice does not by itself pay the
time-integrated budget \eqref{eq:gt-backward-gap}.
